\section{Discussion}\label{sec:discussion}
\paragraph{Design implications}
Four lessons follow for executed cover. First, choose the dispatcher before the cover content: cover that does not pass through a model wait carries a timing signature that removal attacks exploit. Second, matching marginal gaps is insufficient; achieved timing and chronological context remain informative, so a schedule must be evaluated on the records an observer actually sees. Third, planner cover is limited by $P^*(k)$, which falls slowly with more shadows and their linear token cost, whereas drawing shadow labels without replacement lowers it to $1/(k+1)$ at no extra cost; this is the cheapest improvement our analysis predicts and should be measured next. Fourth, isolation is necessary but not sufficient: hidden isolated cover still changes planning paths, and only a state check against declared invariants detects collateral writes.

\paragraph{Threats to validity}\label{sec:limits}
All runs are on one host over loopback, with synthetic tools, two fixed 8B models, and fixtures and phrasings reused across splits and studies. Loopback timing is cleaner than a wide-area path, so the timing attacks may be optimistic for a remote observer; network jitter and loss, cross-host deployment and unseen phrasings are untested. The recipient observer of Table~\ref{tab:cover-necessity}, which also reads decoy arguments, is not measured, and the exact real-chain test for shadow requires per-fixture predictions not reported here. Same-service calls serialize; unknown-task discovery and unrestricted multiplexing are not evaluated. The decision tolerances were set after the measurements, and the recognition intervals in Table~\ref{tab:decision} treat a suite maximum as a single attack; they ignore selection and are therefore too narrow. The record-only suite lacks the adaptive decoders used with call windows, so its counts are not comparable for history-matched cover. Persistent caches and phase order affect costs. The gap of Section~\ref{sec:gap} concerns four settings from two sources, one of them a preprint; twelve of the fifteen mapped studies are preprints.

\paragraph{Open measurements}
The results call for five measurements: a shadow sweep over $k$ and over iid versus distinct label sampling, with the exact real-chain test; context-removal and joint decoders on the record-only representation; the recipient view, using the per-service request journals; jitter, loss and cross-host deployments with unseen phrasings; and equal-budget comparisons of controller and planner cover under production state contracts with legitimate concurrent updates. Supplement S8 specifies further open questions on feedback attribution and task discovery.

\section{Related Work}\label{sec:related}
\paragraph{Agent side channels}
Section~\ref{sec:background} and Appendix~\ref{app:map} cover attacks on Agent traffic. Closest to our setting, WRA evaluates research decoys and destination visibility against a transit observer~\cite{jeong2026wra}, and the communication-graph study evaluates metadata transformations and live OHTTP~\cite{dangol2026workflow}; neither reports matched task utility for executed cover. Serving-side work shows that token streams, cache sharing and token counts leak through model interfaces~\cite{weiss2024prompt,whisperleak2025,earlybird2025,timewilltell2024,carlini2024timing,wu2026cacheprune}.

\paragraph{Decoy-query obfuscation}
Private web search hides a user's queries among machine-generated decoys. TrackMeNot issues periodic decoy queries~\cite{howe2008trackmenot}, and GooPIR submits each query together with $k-1$ decoy keywords to obtain $k$-fold deniability~\cite{domingoferrer2009goopir}. Peddinti and Saxena showed that classifiers separate TrackMeNot decoys from real queries~\cite{peddinti2010trackmenot}, Balsa et al.\ systematized such systems and their failure to provide the deniability they claim~\cite{balsa2012obpws}, and Gervais et al.\ quantified the residual privacy of obfuscated search~\cite{gervais2014websearch}. Our floor $P^*(k)$ is this deniability argument applied to complete Agent executions, and our dispatcher result parallels their lesson that decoys leak through the process that generates them. Two things are new in the Agent setting: decoys perform business work that must satisfy a state contract, and their generation is coupled to model inference, so a controller that bypasses inference creates a timing signature that query generators do not have.

\paragraph{Traffic-analysis defenses}
WF defenses pad, pace and add dummy traffic, and their evaluation methodology, including adapted attacks and matched overheads, is well developed~\cite{juarez2014critical,mathews2023critical,dyer2012peekaboo,wang2017walkietalkie,gong2020front,juarez2016wtfpad,shmatikov2006adaptive,cai2014tamaraw}. Smart-home traffic shaping discards cover at a tunnel endpoint~\cite{apthorpe2019stp}. We carry their evaluation obligations over to cover that executes, and add business state and the dispatcher.

\paragraph{Agent security}
Information-flow and privilege controls such as CaMeL, FIDES and Progent restrict what tools an Agent may call and what information flows between them~\cite{debenedetti2025camel,costa2025fides,shi2025progent}; CaMeL explicitly notes tool-call and timing side channels. SPILLage studies behavioral oversharing by browser Agents~\cite{roh2026spillage}. $\tau$-bench and AgentDojo evaluate environment outcomes~\cite{yao2024taubench,debenedetti2024agentdojo}; our state endpoint follows their practice and binds it to traffic measurements. Surveys organize Agent threats by lifecycle and attack surface~\cite{yao2024survey,shahriar2025survey,ferrag2025survey}.
